\documentclass[10pt]{article}

\usepackage[margin=1in]{geometry}
\usepackage{amsmath}
\usepackage{amssymb}
\usepackage{booktabs}
\usepackage{longtable}
\usepackage{xcolor}
\usepackage{enumitem}
\usepackage[hidelinks]{hyperref}
\usepackage{xurl}
\usepackage{microtype}

\hypersetup{
  colorlinks=true,
  linkcolor=black,
  urlcolor=blue!55!black,
  citecolor=black,
  pdftitle={mirror-pool: behavioral anonymity and confidential value on Solana, verified on-chain},
  pdfauthor={mirror-pool}
}

% code style helpers (identifiers are typeset in a monospace face)
\newcommand{\code}[1]{\texttt{#1}}

\setlength{\parskip}{0.35em}
\setlength{\parindent}{0pt}
\setlength{\emergencystretch}{2.5em}

\title{\textbf{mirror-pool}\\[0.3em]
\large Behavioral anonymity and confidential value on Solana,\\ verified on-chain}

\author{A Superteam Brasil ``Privacy-Through-Noise'' submission}
\date{Design paper, v1}

\begin{document}

\maketitle

\begin{abstract}
\noindent
On-chain behavior is legible. Even when a wallet reveals nothing directly, the
\emph{shape} of what it does (when it acts, how much it moves, who pays its gas,
what software built its transaction, where its funds came from) is enough for
commercial chain-analysis to cluster it, for copy-traders to shadow it, and for
funding provenance to de-anonymize it. mirror-pool is an anonymity system over
on-chain \emph{behavior} rather than over funds, composed of two paths with two
deliberately different strengths and an optional value layer. A
\textbf{synchronized-crowd path} batches $N$ identical actions into one epoch that
settles atomically on a single block timestamp, paid by a rotating gasless relay,
so the strongest empirical attack on mixer-style pools (FIFO temporal matching)
has nothing left to order; each participant still signs their own action, so what
this path destroys is the per-actor \emph{signal} (timing, ordering, size, gas
payer, fingerprint, parseable intent) and not the on-chain signer, which is
collective intent-deniability rather than who-initiated unlinkability. A
\textbf{zero-knowledge opt-in path} is the one that hides WHO: it proves membership
on-chain via Groth16 and settles with no participant signature at settle, so no
wallet is linked to the settled instance, within the set its own public inputs
leave standing (Section~\ref{sec:limits}). An optional
\textbf{confidential-value layer} adds a 2-in/2-out Tornado-Nova JoinSplit that
hides HOW MUCH inside encrypted notes. Running the ZK path with the value layer
composes the two privacy axes; both are proven on-chain on public Solana devnet. We
report a falsifiable anonymity metric rather than a nominal count: an adversarial
harness measures attacker advantage over the $1/k$ guess, and an
information-theoretic effective-$k$ model (Serjantov-Danezis $2^{H(p)}$ plus
min-entropy) shows that a naive pool advertising $k=32$ delivers an effective set
of about $7.5$ (worst case $1$) under funding-provenance partitioning. We answer
that channel with a mechanism rather than an assumption - commit wallets funded by
unshielding from the value pool in denominated, batched funding rounds - and then
measure what it leaves: $24.25$ of $32$ ($75.8\%$) from the rules the protocol can
enforce, or $28.80$ ($90.0\%$) if participants also voluntarily dwell two rounds,
both at full adoption, with the residual published rather than rounded away. That
funding leg is library code plus a CLI that emits a request: it is not wired into a
running service and has never been soaked, and we label it as such. Everything else
is deployed and browser-verifiable on public devnet (program
\code{Eez\-Wd\-Frm\-HtR2\-PCuuc\-Uruvkgy\-B9HW3\-w2Ksk\-ZNeY\-mX\-szBq}),
with two live soaks passing $17/17$ behavioral and $25/25$ confidential on-chain
assertions, backed by $262$ host tests and $99$ on-chain program tests.
\end{abstract}

\section{Introduction}

A classic mixer breaks the link between a deposit and a withdrawal of
\emph{value}. mirror-pool attacks a different property: the behavioral legibility
of an on-chain \emph{action}. $N$ participants voluntarily pool the same action
into one synchronized epoch; an observer sees $N$ identical actions land together
and cannot extract the per-initiator signal that chain-analysis tools depend on:
timing, ordering, amount, gas payer, and wallet fingerprint. The anonymity set is
the set of real participants in the round, exactly as a mixer's anonymity set is
the set of deposits of one denomination. The funds are visible, the action is
visible, the aggregate is visible.

What is destroyed depends on the path, and the two are not interchangeable. On the
\textbf{crowd path} each participant signs their own action inside the shared
settlement transaction, so the action remains attributable to that wallet; what is
destroyed is every per-actor signal an attacker would use to single one
participant out, which defeats copy-trading and FIFO temporal matching and gives
the round collective intent-deniability. On the \textbf{ZK opt-in path} the
settled output is released to a fresh address with no participant signature,
behind an on-chain membership proof, and there the initiator of a given instance
genuinely is unattributable. We state this distinction up front because
overstating it is precisely the failure mode this paper criticizes in prior work.

This is behavioral obscurity for the ``privacy-through-noise'' theme: the goal is
to make on-chain \emph{behavior} hard for automated chain-analysis to read. An
optional confidential-value layer then goes past that theme and also hides
\emph{how much}: value lives in encrypted notes and moves through a
Tornado-Nova-style zero-knowledge JoinSplit, so a private transfer reveals
neither initiator nor amount. Composing the two axes is the point, not scope
creep: a deployment that runs the ZK path together with the value layer hides who
acted and how much moved at once.

\paragraph{Contributions.}
\begin{itemize}[leftmargin=1.4em,itemsep=0.15em]
  \item \textbf{Two peer behavioral paths} over one accumulator, one epoch
        clock, one $k$-floor, and one anti-Sybil economy: a synchronized-crowd
        path (behavioral intent-deniability against copy-trading and per-actor
        signal extraction) and a ZK-deniable opt-in path (cryptographic
        who-initiated unlinkability via an on-chain Groth16 membership proof).
  \item \textbf{An optional confidential-value layer}: a clean-room 2-in/2-out
        JoinSplit (\code{shield} / \code{transfer} / \code{unshield}) with value
        conservation, range proofs, \code{extDataHash} binding, and ECIES notes,
        composable with the ZK path to hide both axes at once.
  \item \textbf{A funding-provenance mechanism, and its measured residual}: the
        commit wallet is funded by unshielding from the value pool in
        denominated, batched funding rounds, so no public edge links a main
        wallet to a commit wallet; the effective-$k$ model derives the resulting
        anonymity from that mechanism's rules instead of assuming it, and
        publishes the leak that survives. This leg is a library plus a CLI that
        emits a funding request; the ingestion wiring and a live soak are not
        done, so it is a designed and measured mechanism, not a running one.
  \item \textbf{A falsifiable anonymity story}: an adversarial A/B harness that
        \emph{measures} FIFO/amount/fingerprint/gas-payer attacker advantage,
        plus an information-theoretic effective-$k$ model on the
        advertised-vs-effective axis the empirical mixer literature made
        memorable.
  \item \textbf{Trustless on-chain verification}: pure-Rust proving (no Node) and
        on-chain Groth16 verification via the BN254 pairing syscalls, so
        settlement checks a proof against a vendored verifying key rather than
        trusting a committee, all deployed and browser-verifiable on public
        devnet.
\end{itemize}

Every empirical claim in this paper is downstream of an executable artifact in
the repository (a harness, a test, or a live soak), and Appendix~\ref{app:repro}
gives the exact command for each. Where the harness disagrees with the prose, the
harness is the source of truth.

\section{Problem and threat model}
\label{sec:threat}

Threat models written for EVM chains under-count Solana's observability, so we
enumerate it first. Solana surfaces \emph{more} structured per-transaction data
than Ethereum, which makes several attacks that are approximate on EVM
\emph{exact} here.

\begin{table}[h]
\centering
\small
\begin{tabular}{@{}p{0.30\linewidth}p{0.62\linewidth}@{}}
\toprule
\textbf{Observable} & \textbf{Why it matters} \\
\midrule
Balance deltas, first-class & \code{pre/postBalances} and \code{pre/postTokenBalances} are transaction metadata; exact per-account deltas are one RPC call away. Amount-matching is cheap and exact; a $0$ pre-balance flags a fresh account. \\
Deterministic ATAs & An Associated Token Account is a pure function of $(\text{owner},\text{mint},\text{program})$; any value landing in a canonical ATA re-links to its owner by table lookup. \\
Signers and fee payer & All signers are in the header; the fee payer is the first signer. Self-paid execution is self-attribution; co-signing merges clusters. \\
Compute-budget / fee fingerprints & \code{SetComputeUnitLimit}/\code{SetComputeUnitPrice}, tx version, instruction ordering, and ALT choice are visible and characteristic of wallet software. \\
Timing despite Gulf Stream & No lingering mempool, but Geyser/Yellowstone/shred observers see transactions pre-block ($<$100ms); post-block every tx has an exact slot and intra-slot position. \\
On-chain risk oracles & Programs can gate on a wallet's risk score at execution time, so a design that leaks the initiator leaks it to counterparty programs too. \\
\bottomrule
\end{tabular}
\caption{What an observer actually sees on Solana. The one thing it lacks (a
durable public mempool) buys a naive design nothing: the parseable intent is
still there, observed a few milliseconds later.}
\end{table}

\paragraph{Adversaries.} We model four, all assumed to run full archival
indexing and none assumed weak.
\begin{description}[leftmargin=1.4em,itemsep=0.12em]
  \item[A1 Chain-analysis clustering.] Account roll-up, ATA reversal, co-signer
        merging, and (the strongest anchor) funding-source anchoring: a fresh
        wallet funded from a KYC-clustered source is re-linked immediately.
  \item[A2 Copy-traders / shadowers.] Parse swap intent (mint pair, amount) from
        pre-block validator streams and replay it within the slot. The lesson:
        \emph{timing obfuscation alone loses}; the per-wallet parseable signal
        must be removed, not delayed.
  \item[A3 The pool operator.] Infrastructure, not a trusted party; modeled
        honest-but-curious to adversarial. It can pad, censor, or reorder, but by
        construction it \emph{cannot} substitute or retarget a committed action,
        nor make a participant act twice (epoch-scoped nullifier PDAs).
  \item[A4 Sybil attacker.] Fills epochs with wallets it controls; if an epoch is
        $1$ honest plus $k{-}1$ attacker wallets, elimination de-anonymizes the
        honest one. This is how small mixers die, and it is why we never count
        nominal participants (Section~\ref{sec:anon}).
\end{description}

\paragraph{The real attacks (with the numbers that justify each defense).} We
target the attacks the empirical literature identifies as the actual killers, not
theoretical ones:
\begin{itemize}[leftmargin=1.4em,itemsep=0.12em]
  \item \textbf{FIFO temporal matching}: linking each deposit to the earliest
        plausible later withdrawal linked up to \textbf{49\%} of deposits on
        small Tornado pools and beat every other heuristic by 15 to 22
        percentage points~\cite{tornado}. It is the single strongest known attack
        on mixer-style systems.
  \item \textbf{Amount matching}: variable and round-number amounts leak a large
        fraction of a claimed set. Wang et al. measured \textbf{27.34\%}
        (Ethereum) and \textbf{46.02\%} (BSC) anonymity-set reduction from
        composable heuristics~\cite{wang}.
  \item \textbf{Wallet fingerprinting}: matching gas/priority-fee settings across
        entry and exit cut Tornado's effective set by \textbf{37\%} on its
        own~\cite{tornado}.
  \item \textbf{Self-paid gas} and \textbf{funding provenance}: not using a
        relayer, and funding from a clustered source, were both dominant
        de-anonymization vectors~\cite{tornado}.
\end{itemize}

\section{Design: behavioral anonymity}
\label{sec:design}

mirror-pool delivers two settlement paths that share one Poseidon accumulator,
one slot-derived epoch clock, one on-chain $k$-floor, and one anti-Sybil economy.
They are peers, not versions; a pool can run either or both.

\subsection{Protocol flow (commit, shared epoch, settle)}

A participant computes a commitment $C = \text{Poseidon}(\text{secret},
\text{actionHash}, \text{epoch})$ off-chain and posts only $C$ on-chain during the
round's commit phase; the action and secret stay client-side. Every commit that
lands in the same slot window belongs to one shared epoch, derived by all
participants from the slot clock with zero coordination
($\text{epoch} = \text{slot} / \text{epoch\_slots}$). Once the window closes at
$\text{settle\_slot} = (\text{epoch}{+}1)\cdot\text{epoch\_slots}$, the off-chain
gasless coordinator checks the real $k$-floor and submits one settlement.

\paragraph{Crowd path (\code{Commit} / \code{SettleEpoch}).} $N$ participants each
sign their own identical action. The coordinator composes them into a single
versioned transaction that settles the whole epoch on one block timestamp:
\begin{center}
\small
\code{ComputeBudget(limit)} $+$ \code{ComputeBudget(price)} $+$
\code{SettleEpoch\{epoch,[nf..]\}} $+$ participant\_0..participant\_{$k{-}1$} actions
\end{center}
Because every action in the epoch shares one timestamp, one fixed
\code{ActionClass}/\code{SizeBucket} shape, and one relay-paid envelope, timing,
amount, gas payer, and wallet fingerprint carry \emph{no bits} about which
participant is worth copying or attributing a strategy to. This is collective
intent-deniability: the action still executes from the participant's named wallet
(so it stays attributable to that wallet), but it is indistinguishable
\emph{as a signal} from the rest of the crowd. \code{SettleEpoch} performs the
anonymity bookkeeping (batching, $k$-floor, per-nullifier anti-replay, authority,
fail-closed parsing); it does not cryptographically bind each nullifier to a
distinct prior commitment, and the paper says so plainly. That binding is exactly
what the ZK path adds.

\paragraph{ZK opt-in path (\code{CommitDeposit} / \code{SettleZk}).} A participant
escrows the action input at commit time, binding a \emph{fresh} output through
$\text{actionHash} = \text{Poseidon}(\text{recipientHi128},
\text{recipientLo128}, \text{amount})$. At settlement the relay verifies a
Groth16 membership proof on-chain that an output corresponds to \emph{some}
member of the accumulator without revealing which, checks the \code{actionHash}
binding and nullifier, and releases the escrow to the bound address with no
participant signature. This provides cryptographic who-initiated unlinkability:
the deposit is visible, no participant signs at settle, and the coordinator is not
in the trust base for attribution (it settles for \emph{some} member without
learning which). The set that unlinkability is over is narrower than the proof
alone: \code{SettleZk} publishes the epoch, so the covering set is the window's ZK
deposits (the amount is public too, but a pool admits exactly one ZK size, so it
separates nothing), and the program enforces no floor on it and no recipient
freshness. Section~\ref{sec:limits} states each of those and why neither can be a
settle-time check on a path whose escrow has no refund.

\subsection{Why the design defeats the named attacks}

\begin{table}[h]
\centering
\small
\begin{tabular}{@{}p{0.30\linewidth}p{0.62\linewidth}@{}}
\toprule
\textbf{Attack} & \textbf{Mechanism that defeats it} \\
\midrule
FIFO temporal matching & Shared-epoch batching: all actions settle on \emph{one} block timestamp, so there is no per-participant settlement time to sort. Per-actor random delay does \emph{not} achieve this (heavy-tailed timing survives); removing the ordering signal does. This mirrors Penumbra batch clearing~\cite{penumbra}. \\
Amount matching & Fixed size buckets baked into the type system: \code{SizeBucket} is a field of \code{ActionClass}, one anonymity set per class, and the bucket is bound into the commitment, so a participant cannot commit one bucket and settle another. \\
Wallet fingerprinting & The coordinator submits every epoch with one pool-wide profile (identical CU limit, priority fee, tx version, account ordering, one shared ALT). There is no per-participant transaction to fingerprint. \\
Self-paid gas & Settlement is gasless for participants; a \emph{rotating} fee-payer set with independent funding histories signs and funds it, so no acting wallet self-attributes and no single relay becomes a consolidation cluster node. \\
Copy-trading / shadowing & Before settlement the only thing on the wire is a commitment hash; at settlement $N$ identical actions execute at once with no per-actor signal, so there is no individual to isolate and shadow. \\
Common funding source & The commit wallet is funded by \emph{unshielding} from the confidential-value pool (\code{fund-commit}), released in denominated, batched funding rounds with a minimum-size floor: the vault is the on-chain sender and the relay the only signer, so no public edge links a main wallet to a commit wallet. This is the one defense with a \emph{measured residual} rather than a clean close (Section~\ref{sec:anon}): the boundary amounts and slots stay public, leaving a matching problem worth $24\%$ of nominal $k$ on the enforced rules, or $10\%$ if participants also dwell two rounds. It is also the one row that is \emph{not wired}: the CLI emits the request, the batcher is a library type, and nothing joins them. \\
\bottomrule
\end{tabular}
\caption{Each defense maps to a concrete mechanism in the codebase, and each is
measured (not asserted) by \code{crates/mirror-harness} (Section~\ref{sec:anon}).
The last row is the exception on both counts: it is a designed mechanism whose
residual is modelled, and it is not connected end to end.}
\end{table}

\subsection{On-chain account and instruction model}

Every account is a PDA whose address is a pure function of the pool and the
epoch/nullifier/participant bytes; the program re-derives the expected address and
rejects any mismatch (\code{InvalidPda}). Every layout is explicit byte offsets
with bounds-checked little-endian accessors: no \code{unsafe}, no transmutes, no
\code{repr(C)} casts on untrusted bytes. Parsing is fail-closed everywhere (wrong
length, unknown tag, out-of-range count, or trailing bytes are rejected before any
account is touched), and the wire layout is defined once and mirrored byte-for-byte
between the host builder and the SBF parser with compile-time size asserts on both.

\begin{table}[h]
\centering
\small
\begin{tabular}{@{}cll p{0.44\linewidth}@{}}
\toprule
\textbf{tag} & \textbf{instruction} & \textbf{who submits} & \textbf{effect} \\
\midrule
0 & \code{InitPool} & operator & fix \code{epoch\_slots}, \code{k\_floor}, \code{entry\_fee}, \code{reward\_bps}, authority (all immutable) \\
1 & \code{Commit} & participant & append one 32-byte leaf, bump commit count, collect fee, accrue dwell \\
2 & \code{SettleEpoch} & rotating relay & enforce authority $+$ $k$-floor, one Nullifier PDA per spend, mark settled \\
3 & \code{CommitDeposit} & participant & escrow \code{amount}, append leaf whose \code{actionHash} binds (recipient, amount) \\
4 & \code{SettleZk} & rotating relay & verify Groth16 membership on-chain, release escrow to the bound recipient \\
5 & \code{ClaimReward} & participant & dwell-proportional, drain-safe reward-pool payout \\
6 & \code{InitValuePool} & operator & create the confidential ValuePool $+$ vault; fix authority, fee, denomination \\
7 & \code{Transact} & rotating relay & verify one 2-in/2-out JoinSplit on-chain; shield / transfer / unshield \\
\bottomrule
\end{tabular}
\caption{The instruction set. The first byte is the discriminator; the four keys
of the account model are the Pool, Epoch, Nullifier, and Dwell PDAs.}
\end{table}

The Pool PDA (1780 bytes) holds one immutable config plus an inline depth-20
Poseidon frontier accumulator (up to about $1.05$M leaves, so an append never
touches a second account), a $32$-root history ring (a membership proof is made
against a root \emph{snapshot}, so settle must accept any recent root), and the
strictly-additive incentive counters (offsets $0\ldots1762$ are byte-identical to
the pre-incentive layout). Anti-replay is one PDA per $(\text{pool},
\text{epoch}, \text{nullifier})$: its existence marks the nullifier spent, and
scoping by epoch means the same secret yields a fresh, unlinkable nullifier in a
later epoch. Immutability is a security property: a mutable $k$-floor would let an
operator lower the floor right before a targeted epoch and shrink the set on
demand.

\paragraph{Solana specifics.} A settlement uses a v0 transaction with an Address
Lookup Table to keep shared non-signer accounts (program ids, Pool PDA, clock,
shared sink) out of the static list, while per-settlement accounts (Epoch PDA,
Nullifier PDAs) and all signers stay static (Solana forbids loading a signer
through an ALT). The message must fit the 1232-byte packet limit, which caps the
fixed-shape \code{PlainTransfer} baseline at $4$ participants per settlement
transaction ($5$ serialize to $1234$ bytes); \code{plan\_settlements} chunks a
larger epoch into transaction-sized groups with disjoint nullifier subsets whose
union is the whole epoch. Because Jito bundles cap at $5$ transactions, $N$-party
atomicity is program-side: \code{SettleEpoch} settles all intents in one
instruction, so a no-show simply forfeits its slot and fee without stalling the
epoch.

\section{Confidential-value layer}
\label{sec:value}

Everything above works on the behavioral axis while leaving every amount public.
The confidential-value layer is the optional complement: a Tornado-Nova-style
shielded pool that hides how much moves. It is a \emph{separate} subsystem (its
own account, accumulator, and two instructions), so a deployment can run the
behavioral pool, the confidential pool, or both. A confidential transfer settles
through the same gasless relay with no user signature and carries
$\text{publicAmount} = 0$, so an observer learns neither which wallet initiated
\emph{that transfer} nor how much it moved. Composing the ``who $+$ how much''
axes onto a single pooled action means running this layer with the ZK opt-in path;
on the crowd path the amounts are hidden but the action still carries the
participant's own signature. The layer is a clean-room implementation of the public, open-source
Tornado-Nova transaction circuit (the standard Poseidon JoinSplit), cited as
prior art.

\paragraph{The value-note (UTXO) scheme.} A value note is
$\{\text{amount}, \text{public\_key}, \text{blinding}\}$, owned by a value keypair
distinct from any Solana wallet:
\begin{align*}
\text{public\_key} &= \text{Poseidon}(\text{private\_key}) \\
\text{commitment}  &= \text{Poseidon}(\text{amount}, \text{public\_key}, \text{blinding}) \quad(\text{Merkle leaf}) \\
\text{signature}   &= \text{Poseidon}(\text{private\_key}, \text{commitment}, \text{leafIndex}) \\
\text{nullifier}   &= \text{Poseidon}(\text{commitment}, \text{leafIndex}, \text{signature})
\end{align*}
The nullifier is bound to both the owner (through \code{signature}, which needs
the private key) and the leaf position, so the same note at a different index
yields a different nullifier and a double-spend collides. A dummy input has
$\text{amount}=0$; the circuit skips its Merkle-membership check, which is what
lets a shield spend two dummy inputs.

\paragraph{One universal statement.} A single 2-in/2-out JoinSplit
(\code{Transact}, tag 7), distinguished only by the signed \code{publicAmount},
covers all three operations:
\[
\text{shield}: \text{publicAmount}=v; \qquad
\text{unshield}: \text{publicAmount}=r-v; \qquad
\text{transfer}: \text{publicAmount}=0.
\]
The two ranges (shield in $[0,2^{248})$, unshield in $(r-2^{248}, r)$) are
disjoint, so the on-chain decoder recovers the sign unambiguously.
$\text{extDataHash} = \text{keccak256}(\text{recipient} \Vert \text{relayer}
\Vert \text{fee\_be} \Vert \text{enc0} \Vert \text{enc1}) \bmod r$ is recomputed
on-chain and required to match, so the relay cannot retarget the recipient, change
the fee, or swap the encrypted payloads. The circuit
(\code{circuits/transaction.circom}, depth 20, \textbf{27278 R1CS constraints},
7 public inputs, 56 private inputs) enforces, per input, the key / commitment /
signature / nullifier relations plus Merkle inclusion when $\text{amount}\neq0$;
per output, the commitment and a 248-bit range bind; and globally, value
conservation $\sum \text{inAmount} + \text{publicAmount} = \sum \text{outAmount}$,
distinct input nullifiers, and tamper-evidence of \code{extDataHash}. The
on-chain handler runs its checks in a fixed fail-closed order: authority, recent
root, denomination (if pinned, checked before the expensive work), \code{extDataHash},
create each input nullifier PDA (an all-zero dummy sentinel is skipped), verify
the proof, append both output commitments, move lamports per \code{publicAmount}
keeping the vault rent-exempt, and emit the encrypted blobs as return data.

\paragraph{Encrypted notes and fixed-denomination mode.} Each output commitment is
opaque on-chain, so the sender posts an ECIES blob (X25519 ECDH $\to$
HKDF-SHA256 $\to$ ChaCha20-Poly1305; a fresh ephemeral key per message makes the
AEAD $(\text{key}, \text{nonce})$ pair unique by construction). The plaintext is
the 40 bytes the recipient does not already know
($\text{amount}(8,\text{BE}) \Vert \text{blinding}(32)$), and the on-chain blob is
a fixed 100 bytes ($\text{ephemeral\_pub}(32) \Vert \text{nonce}(12) \Vert
\text{ct+tag}(56)$). A recipient address splits a BN254 \emph{value} key (spend
authority, bound by the commitment) from an X25519 \emph{viewing} key (encrypt and
discover only), so a viewing key can go to an auditor without granting spend. The
recipient trial-decrypts every blob (\code{scan}), keeps hits, and rebuilds each
note to locate and later spend it. Setting a fixed \code{denomination} pins every
public crossing to one magnitude (\code{DenominationMismatch} otherwise): the
value analog of fixed size buckets, giving amount k-anonymity at the boundary.

This layer is an intentional composition, not scope creep. Its guarantee is
stated precisely: it hides in-pool amounts and, for internal transfers and
unshields, the initiator; it does \emph{not} hide the public shield-in and
unshield-out magnitudes or the pool TVL (a public on-chain balance), and the
verifying key it is deployed with came from a dev/test setup
(Section~\ref{sec:limits}).

\section{Anonymity analysis}
\label{sec:anon}

Overstating an anonymity set is the documented failure mode of prior mixers, so
mirror-pool reports the \emph{real} set. \code{KAnon\{nominal, excluded\}} with
$\text{real\_k} = \text{nominal} - \text{excluded}$ separates the nominal commit
count from the real set (nominal minus operator-owned and detected Sybil decoys),
and an epoch settles only when $\text{real\_k} \geq \text{k\_floor}$. Anonymity is
$1/\text{real\_k}$, never $1/\text{nominal}$. We back this with two measurements.

\subsection{A/B attack harness (advantage over the $1/k$ guess)}

\code{crates/mirror-harness} implements four attacks over deterministic synthetic
populations ($n=2048$ per scenario, fixed seed \code{0x4d4952524f52},
50/50 train/test split) and reports attacker advantage
$\text{Adv} = P[\text{correct attribution}] - 1/\text{real\_k}$ for a
\emph{Baseline} (per-actor delay, self-paid gas, variable amounts, distinct
funding roots) versus \emph{MirrorPool} (shared-epoch batch, rotating gasless
relay, fixed bucket). A positive measured advantage is a broken claim, full stop.

\begin{table}[h]
\centering
\small
\begin{tabular}{@{}lrrrr@{}}
\toprule
& \multicolumn{2}{c}{\textbf{Baseline}} & \multicolumn{2}{c}{\textbf{MirrorPool}} \\
\cmidrule(lr){2-3}\cmidrule(lr){4-5}
attack ($k=16$, guess $=6.25\%$) & acc & adv & acc & adv \\
\midrule
FifoTemporalMatch  & 83.79\% & $+77.54$pp & 5.96\% & $-0.29$pp \\
AmountMatch        & 91.02\% & $+84.77$pp & 6.25\% & $+0.00$pp \\
WalletFingerprint  & 88.77\% & $+82.52$pp & 6.25\% & $+0.00$pp \\
GasPayerReuse      & 84.86\% & $+78.61$pp & 6.64\% & $+0.39$pp \\
\bottomrule
\end{tabular}
\caption{Held-out attribution accuracy and advantage over random guessing at
$k=16$ (full output at $k=4,8,16$ from \code{cargo run -p mirror-harness}). The
headline: FIFO advantage collapses from $+77.54$pp under per-actor delay to
$-0.29$pp under shared-epoch batching (statistically indistinguishable from $0$),
and the amount and fingerprint channels drop to exactly $0.00$pp. Across $k$, FIFO
advantage moves $+72.46 \to -2.73$pp ($k{=}4$), $+79.20 \to +0.20$pp ($k{=}8$),
$+77.54 \to -0.29$pp ($k{=}16$).}
\end{table}

\subsection{Effective-$k$: advertised $k$ is not effective $k$}

A pool that advertises $k$ participants does not provide $k$-anonymity once an
adversary partitions those participants by information it already holds. Following
the standard information-theoretic metrics, we report the Shannon effective set
$2^{H(p)}$ (Serjantov-Danezis~\cite{serjantov}) and the min-entropy worst-case set
$1/\max_i p_i$ (Diaz et al.~\cite{diaz}) over the distribution $p$ an adversary
assigns after clustering committers. The dominant real leak modeled is
\textbf{funding provenance}: a handful of exchange hot-wallets and faucets fund
most participants (modeled as a Zipf common-funder distribution, about one funder
per four committers), so the $k$ committers collapse into a few equivalence
classes and someone with a unique funder is alone in their class (worst case $1$).
On top of provenance we fold in the same timing, amount, and fingerprint channels
the attack battery models.

\paragraph{The mechanism, not an assumption.} mirror-pool answers this channel on
the funding leg: the participant funds a \emph{fresh} commit wallet by unshielding
from the confidential-value pool (\code{mirror-cli fund-commit}), and a funding-round
batcher (\code{mirror\_coordinator::funding}) holds those withdrawals to a round
boundary and releases them in an arrival-independent order. The on-chain sender is
the pool vault and the relay is the only signature, so the
main-wallet-to-commit-wallet edge is never written. What survives is \emph{not}
nothing: \code{publicAmount} is public on both boundary crossings, so an observer
sees the deposits and the withdrawals and is left with a \textbf{matching problem}.
Our effective-$k$ model derives the MirrorPool provenance classes from exactly that
matching, which is what makes the result a measurement rather than a restatement of
the design intent.

\paragraph{Scope of what follows.} Two caveats bind every number in this section,
and we put them before the table rather than after it. First, this leg is
\textbf{not wired}: \code{fund-commit} proves the unshield and prints the emitted
request, \code{FundingRounds} is a unit-tested batcher type, and nothing in the
repository connects one to the other (the coordinator binary is an in-memory
scheduler demo), so no funding round has ever released a withdrawal. Second, the
model separates what the protocol \emph{enforces} from what it can only
\emph{recommend}: denomination and batching are enforced by the program and the
batcher, whereas \emph{dwell} - how many rounds a participant leaves value shielded
before asking for the withdrawal - has no field in \code{FundingRoundConfig} and no
on-chain check. We therefore report the dwell-$0$ row as the guarantee and the
dwell-$2$ row as what cooperative participants reach. Both assume $100\%$ adoption.

\begin{table}[h]
\centering
\small
\begin{tabular}{@{}lrrrr@{}}
\toprule
\textbf{scenario / funding policy} & \textbf{$k$} & \textbf{$2^{H(p)}$} & \textbf{min-ent.} & \textbf{retained} \\
\midrule
\multicolumn{5}{@{}l}{\emph{(a) funding-provenance channel alone (the marquee axis)}}\\
Baseline: public funding edge      & 32 & 7.51  & 7.51  & 23.5\% \\
MirrorPool: pass-through funding   & 32 & 7.66  & 7.51  & 23.9\% \\
MirrorPool: denominated only       & 32 & 10.19 & 7.65  & 31.8\% \\
MirrorPool: batched rounds only    & 32 & 10.93 & 7.66  & 34.2\% \\
MirrorPool: den.\ $+$ rounds, dwell 0 & 32 & \textbf{24.25} & 18.42 & \textbf{75.8\%} \\
MirrorPool: den.\ $+$ rounds, dwell 2 & 32 & 28.80 & 21.00 & 90.0\% \\
\addlinespace
Baseline: public funding edge      & 16 & 5.98  & 5.98  & 37.4\% \\
MirrorPool: den.\ $+$ rounds, dwell 0 & 16 & 12.22 & 9.22  & 76.4\% \\
MirrorPool: den.\ $+$ rounds, dwell 2 & 16 & 14.30 & 10.77 & 89.4\% \\
Baseline: public funding edge      & 64 & 9.47  & 9.47  & 14.8\% \\
MirrorPool: den.\ $+$ rounds, dwell 0 & 64 & 47.11 & 35.93 & 73.6\% \\
MirrorPool: den.\ $+$ rounds, dwell 2 & 64 & 56.85 & 41.17 & 88.8\% \\
\addlinespace
\multicolumn{5}{@{}l}{\emph{(b) all channels (provenance $+$ timing $+$ amount $+$ fingerprint)}}\\
Baseline: public funding edge      & 16 & 1.01 & 1.00 & 6.3\% \\
Baseline: public funding edge      & 32 & 1.02 & 1.01 & 3.2\% \\
Baseline: public funding edge      & 64 & 1.01 & 1.00 & 1.6\% \\
MirrorPool: den.\ $+$ rounds, dwell 0 & 32 & 24.25 & 18.42 & 75.8\% \\
MirrorPool: den.\ $+$ rounds, dwell 2 & 32 & 28.80 & 21.00 & 90.0\% \\
\bottomrule
\end{tabular}
\caption{A naive pool advertising $k=32$ delivers an effective set of about
$\mathbf{7.5}$ (worst case $\mathbf{1}$) under funding-provenance partitioning,
and about $\mathbf{1.0}$ once every channel composes. Routing the funding leg
through the value pool and then behaving naively (pass-through) buys almost
nothing ($7.66$); denominated, batched funding rounds reach
$\mathbf{24.25}$ of $32$ on the enforced rules alone (dwell $0$), and $28.80$ if
participants also dwell two rounds voluntarily. Every MirrorPool row assumes
$100\%$ adoption. The missing $\mathbf{24.2\%}$ (or $10.0\%$ at dwell $2$) is the
residual amount/timing channel of the public boundary crossings, and the
worst-placed committer sits at $3.00$, not $32$. MirrorPool's all-channel rows
equal its provenance-only rows because settlement exposes no per-actor timing,
amount, or fee shape.}
\end{table}

\paragraph{Why the settlement channels contribute nothing.} A behavioral channel
informs the adversary only if settlement actually exposed it. mirror-pool settles
one shared-epoch batch: every action carries one settle slot, one bucket amount,
and one normalized fee shape, so the timing, amount, and fingerprint channels
carry \emph{zero variance} across the batch and contribute nothing. Funding is
therefore the only channel left, which is why it gets a mechanism, an ablation,
and a published residual instead of a paragraph.

\paragraph{The residual, ablated.} Each mitigation alone is insufficient:
denomination alone still leaves a narrow causal window (only the few deposits that
could have funded a given withdrawal are candidates) and batching alone still
leaves the amount, which identifies the deposit directly; both land near a third
of nominal at $k=32$. Dwell (how many rounds a participant may sit shielded)
trades latency for anonymity: $0/1/2/4/8$ rounds measure
$24.25 / 27.43 / 28.80 / 30.01 / 30.88$, i.e. diminishing returns past the harness
default of $2$. Since dwell is a participant choice the protocol cannot enforce,
$24.25$ is the number the mechanism guarantees and everything above it is what
cooperation adds. Adoption bounds everything: at $90\%/75\%/50\%/25\%$ adoption the
dwell-$2$ configuration measures $24.65 / 20.07 / 13.56 / 9.33$, and as soon as one
committer tops up directly the worst case returns to $1.00$. A protocol that
advertises funding privacy while most of its traffic tops up from main wallets is
advertising a number it does not have.

\paragraph{How hard is the attacker?} The matching can be scored two ways, and we
publish both. The weaker attacker normalizes each withdrawal's candidate deposits
independently; the stronger one solves the whole assignment (each deposit funds one
withdrawal), approximated by Sinkhorn scaling because exact matching marginals are
permanents and $\#P$-hard. Every number above is the \emph{stronger} attacker's.
The gap is largest where the causal window is narrowest (denominated-only at
$k{=}64$: $19.57 \to 17.33$) and small under denominated rounds at dwell $2$
($28.87 \to 28.80$ at $k{=}32$). We note honestly that Sinkhorn is an
approximation, not a bound: in one measured epoch it puts slightly \emph{less} mass
on the truth than the independent reading, and the repository pins that instance in
a test so the claim stays "measured stronger on this grid", not "provably stronger".

\paragraph{Sybil dominance closes the ``excluded $=0$'' gap.} When $75\%$ of the
nominal set is attacker-owned decoy traffic, an adversary that owns the decoys
removes them, so effective-$k$ collapses to exactly $\text{real\_k} =
\text{nominal} - \text{excluded}$: nominal $16$/$32$/$64$ with $12$/$24$/$48$
excluded gives effective-$k$ of $4.00$/$8.00$/$16.00$. This makes \code{KAnon}'s
honest subtraction an information-theoretic result rather than a bare assertion:
an inflated nominal set buys no anonymity against anyone who can identify the
decoys.

These effective-$k$ numbers are a \emph{model} (deterministic synthetic
populations, a modeled Zipf common-funder distribution, a modeled funding trace),
honestly labeled as such; an on-chain settlement-trace loader that computes
effective-$k$ over real funding provenance is a named roadmap extension. What is
\emph{measured} is every number in the tables, and in particular that the
mechanism's own residual is nonzero: the repository pins a test asserting the
default policy lands strictly below nominal, so this project cannot drift back
into reporting a perfect score. \code{docs/EFFECTIVE\_K.md} separates measured
from modelled line by line, including the omissions that favour us.

\paragraph{Anti-Sybil economy.} The metric is only as strong as its Sybil
assumptions, so mirror-pool prices identity: each commit carries a fixed,
non-refundable \code{entry\_fee} (fidelity-bond style), collected on both the
crowd \code{Commit} and the ZK \code{CommitDeposit}, so filling an epoch with
$k{-}1$ Sybils costs $(k{-}1)\times\text{fee}$ per epoch forever (the floor rolls
forward until it is met). A \code{reward\_bps} share accrues to an on-chain reward
pool that funds a \emph{dwell} reward: \code{ClaimReward} pays a drain-safe share
proportional to how many distinct epochs a participant has committed into, so
staying (which thickens future epochs) is what pays, not merely showing up. Dwell
advances only through a real, fee-paying commit, at most once per epoch, so it
cannot be minted for free. The ZK path deliberately ships no identity-linked
claim; its anonymity-preserving equivalent (a dwell/age ZK proof) is designed, not
implemented, and the docs say so.

\section{Implementation and on-chain verification}
\label{sec:impl}

\paragraph{Pure-Rust proving (no Node).} The participant CLI (\code{mirror-cli})
proves in-process in pure Rust by default: \code{ark-circom} loads the same
compiled circom \code{.wasm} witness calculator and the same committed
\code{.zkey} that snarkjs uses, computes the witness from typed Rust inputs, and
\code{ark-groth16} over \code{ark-bn254} produces the proof using the
snarkjs-compatible R1CS-to-QAP reduction, so the proof verifies under the same
committed verifying key the on-chain program embeds. A legacy \code{--use-snarkjs}
shell-out remains as an optional fallback. The CLI never self-submits: settlement
is paid and signed by the rotating gasless relay, so the prove commands
\emph{emit} the settlement instruction bytes for the relay instead of sending
them.

\paragraph{Trustless on-chain Groth16 (a check, not a committee).} Both ZK paths
verify their proof on-chain with \code{groth16-solana} (v0.2.0 byte layout) via
the BN254 (\code{alt\_bn128}) pairing syscalls, against a verifying key vendored
into the program from \code{circuits/artifacts}. \code{proof\_a} is emitted
pre-negated so the on-chain path needs no ark serialization at runtime. This is
the crucial contrast with a trusted-committee design: settlement does not trust an
operator's word that a member is valid; it checks a proof, and \code{verify()}
(not \code{verify\_unchecked()}) also rejects any public input not reduced modulo
the field. The membership circuit (\code{circuits/membership.circom}, depth 20,
\textbf{11522 R1CS constraints}, 4 public inputs, 41 private inputs) enforces leaf
recomputation, Merkle inclusion, and the nullifier relation; public inputs are the
fixed order $[\text{root}, \text{nullifierHash}, \text{actionHash},
\text{epoch}]$. Host and on-chain Poseidon (\code{light-poseidon} and
\code{sol\_poseidon}) are byte-identical and both match the circuit, pinned by a
fixture cross-check test that reproduces the circuit's \code{nullifierHash},
commitment leaf, and Merkle root exactly.

\paragraph{Compute cost.} On-chain verification is cheap enough for
per-membership settlement (about 170K to 500K CU for a 128-byte compressed proof);
Section~\ref{sec:eval} gives the measured CU per flow on devnet.

\paragraph{Coordinator.} \code{mirror-coordinator} maps its privacy jobs onto its
modules: shared-epoch batching and the $k$-floor gate, honest \code{KAnon}
accounting, the rotating gasless relay, and the normalized transaction profile
(\code{cu\_limit} $=400{,}000$, priority fee $=10{,}000$ micro-lamports). It is
generic over a submitter seam, so the identical privacy-critical batching and
floor logic runs against an in-memory submitter in tests and the real submitter in
production, with no validator needed to test the decisions that matter.

\section{Trusted setup: a verifiable multi-party ceremony}
\label{sec:ceremony}

Groth16's price for small, cheap on-chain verification is a per-circuit structured
reference string whose sampling secrets (``toxic waste'') must be destroyed. We
implement a distributable multi-party \textbf{phase-2} ceremony for both circuits
(\code{crates/mirror-ceremony}, driven from \code{mirror-cli ceremony}), in pure
Rust over the same arkworks stack the prover uses.

\paragraph{Phase 1 is imported, not generated.} \code{ceremony start} takes a public
perpetual powers-of-tau file and records its SHA-256, curve, power, and contribution
count in the transcript; \code{ceremony inspect-ptau} parses the file's own
contributions section and prints every contributor's self-declared name, so the list
can be compared against the published record rather than asserted in prose. A file
with fewer than two contributions is refused unless explicitly overridden. The
initial phase-2 key is the output of \code{snarkjs groth16 setup}, which is
deterministic: the same R1CS and the same powers-of-tau reproduce it byte for byte,
so the start of the chain is re-derivable from public inputs.

\paragraph{A contribution moves only $\delta$.} With a fresh secret scalar $s$,
$\delta_{G_1}$ and $\delta_{G_2}$ are multiplied by $s$ and the $\delta$-divided
query vectors $h$ and $l$ by $s^{-1}$; every other element of the key is unchanged.
The final $\delta$ is the product of all contributions, giving the standard
\textbf{1-of-N honest} property: the setup is safe if \emph{at least one} contributor
destroyed their scalar. Each contribution carries a Schnorr proof of knowledge of $s$
with respect to the previous key's $\delta_{G_1}$, whose Fiat--Shamir challenge is
bound to the running transcript hash, the contribution index, the contributor
identifier, and the step's kind and provenance; without extractable knowledge of the
ratio, a contributor could publish a point derived from an earlier step and cancel an
honest contributor's randomness, and without the kind binding a closing beacon could
be relabelled into an extra ``contributor''. Entries are chained with SHA-256 over a
canonical, length-prefixed serialization. A beacon is \emph{final}: the tool refuses
to append after one, and verification rejects a chain that has a step after the
beacon or more than one beacon.

\paragraph{Verification is reproducible by anyone.} Given the transcript, the initial
key, and the final key, \code{ceremony verify} recomputes the whole chain and checks
the hash links, every entry hash, every proof of knowledge, a pairing same-ratio
check $e(\delta_{G_1}^{prev}, \delta_{G_2}^{new}) = e(\delta_{G_1}^{new},
\delta_{G_2}^{prev})$ per step, exact recomputation of beacon steps from their
published source, the endpoint digests, byte-equality of every $\delta$-independent
part of the key, and a batched pairing check that $h$ and $l$ were divided by exactly
the accumulated ratio. It rejects a tampered delta, a forged or replayed proof of
knowledge, a reordered chain, a truncated chain, a step appended after the closing
beacon, and a beacon relabelled as a contributor; each rejection has a test,
including the sophisticated variants where the adversary re-hashes the entire chain
after tampering so only the algebra can object. Truncating \emph{both} the transcript
and the key is a legitimately shorter ceremony that the algebra cannot reject, which
is precisely why the final transcript hash is the value a coordinator publishes.

\paragraph{Counting contributors honestly.} The tool reports an
\emph{independent-contributor} count, not a contribution count. Contributions sharing
an identity, a machine fingerprint, or a proof-of-knowledge nonce are merged;
deterministic (fixed-seed) and beacon steps are never counted at all, because their
scalars are reproducible by anyone. On a run where three self-asserted identities
contributed from one machine, the tool reports \textbf{1 independent contributor from
4 steps}, not 4. We state its limits explicitly: identifiers and fingerprints are
self-asserted, so this is a guard against accidental self-inflation and
\emph{not} a Sybil defence; it can also under-count when environments look alike,
which is the safe direction. One limit is worth stating precisely, because it is
information-theoretic rather than an implementation gap: a scalar carries no evidence
of where it came from, so a verifier who does \emph{not} hold the pre-committed
beacon value cannot distinguish a public beacon scalar from a secret contributor's.
Binding the kind into the proof of knowledge stops any third party from relabelling a
published transcript, and passing the announced beacon value to \code{ceremony verify}
turns the remaining case into a mechanical check; where that value is absent, the
report says the check did not run rather than implying it did.

\paragraph{End-to-end.} \code{ceremony prove-check} runs a ceremony, generates a real
membership proof under the ceremony-produced proving key, and verifies it with the
\emph{exact} on-chain \code{groth16-solana} verifier against the ceremony-exported
verifying key. In the exported key only \code{vk\_delta\_2} differs from the dev key,
as phase-2 theory predicts, and a ceremony-key proof is correspondingly rejected by
the old dev verifying key. \code{snarkjs groth16 verify} accepts the same proof under
the exported \code{verification\_key.json}.

\paragraph{What has not been done.} No production ceremony has been \emph{run}. The
committed and deployed verifying keys still come from the insecure dev setup, and
that remains true until a ceremony is run with external contributors and the program
is redeployed with its key (Section~\ref{sec:limits}).

\section{Evaluation: public devnet deployment}
\label{sec:eval}

Both soak suites were run against \textbf{public Solana devnet}
(\code{api.devnet.solana.com}), a real, shared, public cluster. Every signature
below is a real devnet transaction that resolves on the Solana Explorer; anyone
can click through and verify it. This is devnet, not mainnet-beta: it proves the
program deploys and every flow executes on a live public cluster exactly as it
would on mainnet, without spending mainnet SOL. The SHA-256 of the on-chain
program bytes equals the SHA-256 of the locally built \code{mirror\_pool.so}, so
the deployed program is exactly the source in this repository.

\begin{itemize}[leftmargin=1.4em,itemsep=0.1em]
  \item \textbf{Program id:} \code{EezWdFrmHtR2PCuucUruvkgyB9HW3w2KskZNeYmXszBq}\\
        \href{https://explorer.solana.com/address/EezWdFrmHtR2PCuucUruvkgyB9HW3w2KskZNeYmXszBq?cluster=devnet}{explorer.solana.com/address/Eez\ldots szBq (devnet)}
  \item \textbf{Behavioral crowd settle} (4 identical actions, one atomic tx, one
        timestamp):
        \href{https://explorer.solana.com/tx/4tcgNnhbZe7YKM26F6SnXW1WctASqVEqQ9W4v4NQ85fDfkYe3eq7YqCr4n7bEogm7U5By17Lkc5Tqa7jmZx693NB?cluster=devnet}{tx 4tcg\ldots 93NB}
  \item \textbf{Confidential private transfer} (\code{publicAmount == 0}, relay-only signer):
        \href{https://explorer.solana.com/tx/2r6mo2YrDtjJP8bnVqRJXaVDgVLkxaAmtKAHvgkwYvWb8rchvd63QinFqbmuj5Gk2Se8M9g9ajcCkC5R6QNprtyv?cluster=devnet}{tx 2r6m\ldots rtyv}
\end{itemize}

\subsection{Behavioral soak: crowd $+$ ZK settlement ($17/17$)}

A fresh pool per run (\code{epoch\_slots}$=128$, \code{k\_floor}$=3$,
\code{entry\_fee}$=10^6$ lamports, \code{reward\_bps}$=2500$). Four participants
commit the \emph{same} \code{PlainTransfer} into one shared epoch, settled by one
atomic gasless transaction (the sink is credited $4\times10^7$ lamports, i.e.\
$4\times$ the $10^7$ bucket); a ZK opt-in escrow of $5\times10^7$ is settled by an
on-chain Groth16 \code{SettleZk}. The run makes one ZK deposit, so that settle
exercises the mechanism (proof, binding, nullifier, escrow move) and not
anonymity: its window holds a single deposit. The adversarial cases all
fail closed on-chain: under-floor no-settle (\code{BelowKFloor}, custom error
\code{0x2}), duplicate nullifier (\code{NullifierSpent}, \code{0x3}), re-settle
(\code{EpochAlreadySettled}, \code{0x6}), mismatched recipient
(\code{ActionHashMismatch}, \code{0xe}), and ZK replay (\code{NullifierSpent}).
All $17/17$ assertions passed.

\begin{table}[h]
\centering
\small
\begin{tabular}{@{}lr@{\hskip 2.2em}lr@{}}
\toprule
\textbf{flow} & \textbf{CU} & \textbf{flow} & \textbf{CU} \\
\midrule
\code{init\_pool}          & 21{,}766 & \code{underfloor\_commit}   & 39{,}499 \\
\code{crowd\_commit} (x4)  & 39{,}647 to 41{,}006 & \code{crowd\_settle}        & 21{,}111 \\
\code{zk\_deposit\_commit} & 42{,}447 & \code{zk\_settle} (Groth16) & 102{,}115 \\
\bottomrule
\end{tabular}
\caption{Measured compute units per behavioral flow, finalized devnet
transactions. The on-chain Groth16 membership verification (\code{zk\_settle})
runs in about 102K CU, comfortably inside a transaction's budget alongside the
settlement work.}
\end{table}

\subsection{Confidential-value soak: shield / transfer / unshield ($25/25$)}

Fresh pools per run: a main ValuePool and a fixed-denomination ValuePool
($10^7$ lamports). Amounts: shield $5\times10^7$, hidden transfer $2\times10^7$,
withdraw $2\times10^7$; the relay fee ($5000$ lamports) is bound into ext-data.
The assertions prove the headline claim on-chain: the vault is credited by the
shield deposit and the value root advances with both output commitments inserted;
the private transfer carries \code{publicAmount == 0} (verified as $32$ zero bytes
in the on-chain \code{Transact} data) and moves \emph{no} public lamports, yet
advances the root and spends its input notes; a recipient auto-discovers the
payment note by scanning encrypted blobs; an unshield credits a fresh recipient
and debits the vault by exactly the withdrawn amount; a wrong-denomination deposit
is rejected on-chain (\code{DenominationMismatch}, \code{0x17}) and refused
client-side; a mutated public input is rejected (\code{ProofVerificationFailed},
\code{0xc}); and conservation holds ($\text{vault} = \text{baseline} +
\text{net deposit} - \text{net withdrawal}$). All $25/25$ assertions passed.

\begin{table}[h]
\centering
\small
\begin{tabular}{@{}lr@{\hskip 2.2em}lr@{}}
\toprule
\textbf{flow} & \textbf{CU} & \textbf{flow} & \textbf{CU} \\
\midrule
\code{init\_value\_pool\_main}  & 28{,}964 & \code{transfer} (JoinSplit) & 197{,}616 \\
\code{init\_value\_pool\_denom} & 33{,}503 & \code{unshield} (JoinSplit) & 202{,}843 \\
\code{shield} (JoinSplit)       & 196{,}837 & \code{denom\_shield\_exact}  & 197{,}077 \\
\bottomrule
\end{tabular}
\caption{Measured compute units per confidential flow, finalized devnet
transactions. Each 2-in/2-out JoinSplit (7 public inputs) verifies on-chain in
about 197K to 203K CU.}
\end{table}

\paragraph{Test suite.} Beyond the two live soaks, the repository ships
\textbf{264 host tests} (\code{cargo test --workspace}, covering the Poseidon
scheme cross-check against the circuit, the wire layout, \code{KAnon} accounting,
the coordinator batching, $k$-floor and funding-round logic, ECIES notes, the
trusted-setup ceremony, and the effective-$k$ and funding-model functions; $11$
further tests are environment-gated and skipped by default) and \textbf{105
on-chain program tests} (mollusk integration tests plus in-crate unit tests).
\code{cargo build-sbf} is green and \code{cargo check --workspace} passes.

\section{Honest limitations}
\label{sec:limits}

A privacy tool that is vague about its non-goals is a trap. We list what still
leaks and what we do not claim.

\begin{itemize}[leftmargin=1.4em,itemsep=0.15em]
  \item \textbf{The DEPLOYED trusted setup is dev/test.} Both committed verifying
        keys (membership and the JoinSplit) came from a single-contributor,
        reproducible dev setup whose phase-2 entropy is a hard-coded public string,
        so their toxic waste is public. The keys verify and the program embeds them,
        but the system MUST NOT secure real value until a real ceremony output is
        exported and the program is redeployed with it. A multi-party phase-2
        ceremony is implemented and runnable (Section~\ref{sec:ceremony}); running
        one with external contributors is operational work that has not been done.
        This is a soundness caveat, disclosed here rather than buried.
  \item \textbf{Devnet, not mainnet.} The deployment proves execution on a live
        public cluster; it is not a mainnet-beta production deployment, and the
        system is not audited. Do not rely on it to protect real activity.
  \item \textbf{Anti-Sybil is economic, not a Sybil oracle.} Real $k$-anonymity
        assumes participants are economically distinct. The per-identity entry fee
        raises the price of flooding a round, but same-funding-source Sybils are
        not computable from the 32-byte commit stream and inflate \code{real\_k}.
        The honest worst case: your anonymity is at least the number of
        participants you personally believe are independent, and at most
        \code{real\_k}.
  \item \textbf{Effective-$k$ is a model.} The numbers come from deterministic
        synthetic populations, a modeled Zipf common-funder distribution, and a
        modeled funding trace, not a live pool. A live-provenance loader is a named
        roadmap item.
  \item \textbf{The funding path is not wired and not soaked.} The unshield it is
        built on is soak-proven on devnet ($25/25$), but the layer above it is two
        disconnected halves: \code{fund-commit} proves an unshield to a fresh
        wallet and \emph{prints} the emitted request, and \code{FundingRounds} is a
        unit-tested batcher type that would hold such requests to a round boundary.
        No shipped component ingests one into the other - the coordinator binary is
        an in-memory scheduler demo - so no funding round has ever run, and the two
        live soaks do not exercise this leg. Wiring the ingestion path and then
        soaking a funding round is named future work, not an implied result.
  \item \textbf{Funding provenance is reduced, not closed, and the headline number
        is conditional.} The funding mechanism turns a public edge into a matching
        problem, and the measurement says how much of that matching survives:
        $75.8\%$ of nominal at $k{=}32$ under denominated batched rounds on the
        rules the protocol can actually enforce (a $24.2\%$ residual), and $90.0\%$
        only if every participant also voluntarily dwells two rounds, which nothing
        in the protocol requires. Both figures assume $100\%$ adoption; it drops to
        $42.4\%$ at $50\%$ adoption, and to $23.9\%$ if participants pass value
        straight through a free-amount pool. The mechanism also does not hide that a
        wallet deposited into the value pool, and a participant whose deposit is
        itself the first hop out of a KYC'd exchange is still identified as a pool
        user.
  \item \textbf{The crowd path is behavioral, not cryptographic.} On the crowd
        path the participant signs their own action, so it stays attributable to
        that wallet, and the coordinator learns the participant-to-intent mapping
        while composing settlement (it reveals nothing the signed on-chain action
        did not). Use the ZK opt-in path for cryptographic who-initiated
        unlinkability. This is a disclosed distinction, not a hidden one.
  \item \textbf{Per-epoch, not cross-epoch.} An adversary with a prior over many
        epochs (intersection attacks) can narrow the set across epochs; neither
        path makes a cross-epoch claim, and mitigations are roadmap items.
  \item \textbf{Public boundary and membership.} Membership is public (like
        Tornado deposits); the confidential layer's shield-in and unshield-out
        magnitudes and the pool TVL stay public; a $k$-floor roll-forward is
        observable; and a participant who sweeps a fresh output into a clusterable
        wallet re-links themselves. Compliance and association-set tooling is
        future work.
  \item \textbf{The ZK path's set is per-window and per-amount, and no on-chain
        floor bounds it.} \code{SettleZk} publishes the epoch and the amount, so
        the covering set is that window's ZK deposits of the same amount, and the
        program will settle into a set of one: it reads neither \code{k\_floor}
        nor any Epoch account. This asymmetry with \code{SettleEpoch} is
        deliberate. A crowd epoch below the floor \emph{rolls forward} at no cost
        to anyone, whereas the ZK escrow's only exit is a \code{SettleZk} bound to
        one (recipient, amount, epoch), with no refund and no re-binding, so a
        settle-time floor would convert thin anonymity into permanently stranded
        escrow. Only the secret holder can produce the proof, so no third party
        can force a thin settle, and the floor is enforced at proof time instead
        (the CLI refuses below \code{k\_floor} unless explicitly waived). For the
        same reason recipient \emph{freshness} is a client convention rather than
        a check: a dusted recipient address would otherwise strand the escrow.
        Both properties are pinned by mollusk tests.
  \item \textbf{The ZK escrow shared-pot hole is closed.} It previously was one:
        the leaf is opaque, so nothing on-chain tied a settled amount to what its
        owner escrowed, and because both paths appended the same leaf shape to one
        accumulator, a fee-only crowd \code{Commit} leaf satisfied the membership
        circuit as well as a deposit leaf did. Both halves are fixed: the ZK pool
        carries a fixed denomination set at init, and crowd leaves are
        domain-separated \emph{by the program},
        $\mathit{crowdLeaf} = \mathrm{Poseidon}(\mathit{CROWD\_LEAF\_DOMAIN},
        \mathit{commitment})$. The tag must be applied by the program to the free
        path rather than by the client to the deposit path, since a deposit leaf is
        32 caller-supplied bytes and a client would simply pre-compute the tagged
        value. Two tests pin it: the original attack now failing, and the
        pre-hashing dodge failing. No circuit change was required, so no verifying
        key was regenerated.
\end{itemize}

\section{Related work}
\label{sec:related}

mirror-pool is an independent clean-room implementation built from public
techniques and papers. \textbf{Tornado Cash} established the
commitment/nullifier/Merkle anonymity set for funds; mirror-pool reuses that
accumulator shape but for \emph{initiators} of behavior, and its behavioral core
is the complement of a mixer (amounts public; on the crowd path the initiator
\emph{signal} is unattributable, on the ZK path the initiator itself is).
\textbf{Tornado-Nova} contributed the 2-in/2-out Poseidon JoinSplit that the
confidential-value layer implements clean-room. \textbf{Privacy Pools} added
association-set membership to let honest users dissociate from illicit deposits; a
mirror-pool association-set layer is named future work rather than a current
claim. \textbf{Penumbra}'s batch swaps, where all swaps in a block clear at one
price so per-user ordering is eliminated by construction, are the direct precedent
for shared-epoch settlement: the per-actor signal removed structurally rather than
through added noise~\cite{penumbra}. Penumbra additionally shields the swap itself,
which mirror-pool's crowd path does not; that stronger property is what our ZK
opt-in path buys separately.
\textbf{Token-2022 confidential transfers} hide the transferred amount while
leaving sender, receiver, and mint public, so the entity graph stays intact; they
are the complement to the behavioral core (amounts hidden but behavior legible),
and composing an amount-privacy primitive with a behavior-privacy one is exactly
the ``who + how much'' point.

\paragraph{An honest tradeoff: transparent setups.} mirror-pool uses Groth16,
which needs a per-circuit trusted setup (Section~\ref{sec:ceremony}; the deployed
keys are still from the dev setup). Approaches built on
STARKs or other transparent, post-quantum setups (for example Bulletproofs-style
range proofs or STARK-based systems) avoid a trusted setup entirely, at the cost
of larger proofs or higher verifier work. That is a real tradeoff, not a strict
win for either side: Groth16 buys small, constant-size proofs and cheap on-chain
pairing verification (the $\sim$100K to 200K CU measured in
Section~\ref{sec:eval}) in exchange for a ceremony; a transparent setup removes
the ceremony risk in exchange for proof size or verification cost. mirror-pool's
choice is pragmatic for on-chain Solana verification today, and its trusted-setup
limitation is disclosed rather than obscured.

\section{Conclusion}

mirror-pool composes two privacy axes that the literature usually treats
separately: on-chain \emph{behavior} and \emph{how much} moved (a
confidential-value JoinSplit). On the behavioral axis it ships two paths with two
different strengths, and says which is which: synchronized-crowd batching removes
the per-actor signal while leaving each action attributable to the wallet that
signed it, and the ZK opt-in path, with an on-chain membership proof and no
participant signature at settle, is the one that hides \emph{who} initiated. It
refuses to overstate anonymity, reporting a falsifiable metric backed by an
adversarial harness and an information-theoretic effective-$k$ model whose headline
is given as the number the protocol can enforce rather than the number cooperative
users can reach, and it verifies its zero-knowledge proofs trustlessly on-chain
rather than trusting a committee. Every headline
number in this paper is reproducible from the repository, and the whole system is
deployed and browser-verifiable on public devnet. Appendix~\ref{app:repro} makes
each number checkable.

\appendix
\section{Reproduce every number}
\label{app:repro}

All commands run from the repository root. The host workspace is one Cargo
workspace; the on-chain program is a separate SBF crate excluded from it.

\paragraph{Build.}
\begin{quote}\footnotesize\ttfamily
cargo build --workspace\\
cargo build-sbf --manifest-path programs/mirror-pool/Cargo.toml
\end{quote}

\paragraph{Tests ($220$ host passing plus $7$ environment-gated; $44$ program).}
\begin{quote}\footnotesize\ttfamily
cargo test --workspace\\
cargo test --manifest-path programs/mirror-pool/Cargo.toml
\end{quote}
The 7 gated tests need build artifacts that are not committed (the compiled
circuits, the powers-of-tau, the proving keys); they are \code{\#[ignore]}d and
additionally require \code{MIRROR\_PROVE\_LIVE=1}.

\paragraph{Trusted-setup ceremony (Section~\ref{sec:ceremony}).}
\begin{quote}\footnotesize\ttfamily
D=ceremony/membership\\
mirror-cli ceremony inspect-ptau -{}-ptau circuits/pot16\_final.ptau\\
snarkjs groth16 setup circuits/membership.r1cs \textbackslash\\
\ \ circuits/pot16\_final.ptau circuits/membership\_0000.zkey\\
mirror-cli ceremony start -{}-circuit membership -{}-dir \$D \textbackslash\\
\ \ -{}-r1cs circuits/membership.r1cs -{}-ptau circuits/pot16\_final.ptau \textbackslash\\
\ \ -{}-initial-zkey circuits/membership\_0000.zkey\\
mirror-cli ceremony contribute -{}-dir \$D -{}-id "alice@example.org"\\
mirror-cli ceremony contribute -{}-dir \$D -{}-id "bob@example.net"\\
mirror-cli ceremony beacon -{}-dir \$D -{}-id coordinator \textbackslash\\
\ \ -{}-source-hex \emph{<pre-committed value>} -{}-iterations-exp 16\\
mirror-cli ceremony verify -{}-dir \$D \textbackslash\\
\ \ -{}-r1cs circuits/membership.r1cs -{}-initial-zkey circuits/membership\_0000.zkey \textbackslash\\
\ \ -{}-beacon-source-hex \emph{<pre-committed value>} -{}-beacon-iterations-exp 16\\
mirror-cli ceremony prove-check -{}-dir \$D -{}-out-dir /tmp/cproof
\end{quote}
The last command proves the membership circuit under the ceremony-produced key and
runs the on-chain \code{groth16-solana} verifier over the result; with
\code{-{}-out-dir} the same proof can be re-checked with \code{snarkjs groth16
verify}. The full guide, including what the independent-contributor count can and
cannot detect, is in \code{docs/CEREMONY.md}.

\paragraph{Adversarial harness (the attack table and the effective-$k$ table).}
\begin{quote}\footnotesize\ttfamily
cargo run -p mirror-harness -{}-release
\end{quote}
Prints the Baseline-vs-MirrorPool advantage table (Section~\ref{sec:anon}, FIFO
$+77.54 \to -0.29$pp at $k{=}16$) and the effective-$k$ tables (funding-provenance
alone: $k{=}32 \to 7.51$ for a public funding edge, $7.66$ for pass-through
shielded funding, $24.25$ for denominated batched rounds at dwell $0$ and $28.80$
at the harness default dwell of $2$; the funding-policy and adversary-strength
ablations; the dwell and adoption sweeps; and the Sybil-dominance case collapsing
to \code{real\_k}). Deterministic (seed
\code{0x4d4952524f52}), so any reviewer re-derives them bit-for-bit.

\paragraph{Public devnet soaks (the $17/17$ and $25/25$ assertion tables and the
CU numbers).}
\begin{quote}\footnotesize\ttfamily
PROG=EezWdFrmHtR2PCuucUruvkgyB9HW3w2KskZNeYmXszBq\\
cargo run -p mirror-soak -{}-\ \\
\ \ -{}-rpc-url https://api.devnet.solana.com -{}-program-id \$PROG\ \\
\ \ -{}-epoch-slots 128 -{}-k-floor 3 -{}-zk-amount 50000000\\[0.3em]
cargo run -p mirror-soak -{}-bin mirror-soak-value -{}-\ \\
\ \ -{}-rpc-url https://api.devnet.solana.com -{}-program-id \$PROG\ \\
\ \ -{}-shield-amount 50000000 -{}-transfer-amount 20000000 -{}-denomination 10000000
\end{quote}
The full reproduce recipe (fresh program keypair, single funded master payer,
deploy, and re-run against a fresh devnet or a local mainnet-mirror validator) is
in \code{docs/PROOF.md}, which also lists every finalized transaction signature.

\paragraph{Browser verification (no build required).} Open the program address and
the two representative transactions on the Solana Explorer:
\begin{itemize}[leftmargin=1.4em,itemsep=0.05em]
  \item Program: \url{https://explorer.solana.com/address/EezWdFrmHtR2PCuucUruvkgyB9HW3w2KskZNeYmXszBq?cluster=devnet}
  \item Crowd settle: \url{https://explorer.solana.com/tx/4tcgNnhbZe7YKM26F6SnXW1WctASqVEqQ9W4v4NQ85fDfkYe3eq7YqCr4n7bEogm7U5By17Lkc5Tqa7jmZx693NB?cluster=devnet}
  \item Confidential transfer: \url{https://explorer.solana.com/tx/2r6mo2YrDtjJP8bnVqRJXaVDgVLkxaAmtKAHvgkwYvWb8rchvd63QinFqbmuj5Gk2Se8M9g9ajcCkC5R6QNprtyv?cluster=devnet}
\end{itemize}

\paragraph{Where each claim lives.} The threat model and the per-attack numbers
are in \code{docs/THREAT\_MODEL.md}; the two settlement paths and the account
model in \code{docs/ARCHITECTURE.md}; the effective-$k$ metric and its honest
limits in \code{docs/EFFECTIVE\_K.md}; the incentive economy in
\code{docs/INCENTIVES.md}; the trusted-setup ceremony, how to contribute to one and
how to verify somebody else's, in \code{docs/CEREMONY.md}; and the live results,
signatures, and CU per flow in \code{docs/PROOF.md}.

\begin{thebibliography}{9}
\small
\bibitem{tornado} \emph{Deanonymizing Tornado Cash: an empirical analysis of
mixed flows.} arXiv:2510.09433. FIFO temporal matching links up to 49\% of
deposits on small pools ($+$15 to 22pp over other heuristics); wallet
fingerprinting via gas/priority-fee settings reduces the effective set by 37\%
alone; self-paid (non-relayer) gas is a dominant de-anonymization vector.

\bibitem{wang} Y.\ Wang et al. \emph{On How Zero-Knowledge Proof Blockchain
Mixers Improve, and Worsen User Privacy.} arXiv:2201.09035. 27.34\% (Ethereum) /
46.02\% (BSC) anonymity-set reduction from composable heuristics; anonymity
mining attracts privacy-indifferent users who degrade the set.

\bibitem{serjantov} A.\ Serjantov and G.\ Danezis. \emph{Towards an Information
Theoretic Metric for Anonymity.} PET 2002. Defines the effective anonymity-set
size as $2^{H(p)}$.

\bibitem{diaz} C.\ Diaz, S.\ Seys, J.\ Claessens, B.\ Preneel. \emph{Towards
Measuring Anonymity.} PET 2002. The companion entropy metric and the min-entropy
/ worst-case reading of anonymity as $1/\max_i p_i$.

\bibitem{penumbra} Penumbra batch swaps. All swaps in a block execute at one
clearing price, so per-user ordering and attribution are eliminated by
construction and only the batch aggregate is revealed; the precedent for
shared-epoch settlement.

\bibitem{token2022} Token-2022 Confidential Transfers (Solana Program Library).
Hide the transferred amount while leaving sender, receiver, and mint public; the
complement to mirror-pool's behavioral core.
\end{thebibliography}

\end{document}
